\chapter*{Qué parece valorar el examen disponible}\addcontentsline{toc}{chapter}{Qué parece valorar el examen disponible}
A partir de la corrección ordinaria de enero de 2025 disponible en el material proporcionado y en la página pública de la asignatura en Wuolah, puede afirmarse que el examen disponible exige encadenar equilibrio global, leyes de esfuerzos y análisis de sección. No se infieren frecuencias estadísticas: se utiliza únicamente como referencia de estilo y nivel.\autocite{wuolah-examen-2025,wuolah-ule-page}

% ============================================================

\safechapter{Cómo resolver problemas de leyes de esfuerzos}

\safesection{Algoritmo general}
Utiliza siempre este orden:

\begin{tcolorbox}[
  colback=grisclaro,
  colframe=gris,
  boxrule=0.6pt,
  arc=1.5mm,
  left=6pt,right=6pt,top=6pt,bottom=6pt
]
\centering
\small
\textbf{1. Modelo} $\rightarrow$ \textbf{2. Reacciones} $\rightarrow$ \textbf{3. Tramos} $\rightarrow$ \textbf{4. Cortes}\\[1.5mm]
$\rightarrow$ \textbf{5. $N,V,M$} $\rightarrow$ \textbf{6. Diagramas} $\rightarrow$ \textbf{7. Comprobaciones}
\end{tcolorbox}

\safesection{Paso 1: dibuja el modelo}
Marca claramente $A$, $B$, $C$, etc., además de longitudes, cargas, apoyos y ejes.

\safesection{Paso 2: calcula reacciones}
Sustituye apoyos por reacciones. Después:
\[
\sum F_x=0,\qquad\sum F_y=0,\qquad\sum M=0.
\]

\safesection{Paso 3: divide en tramos}
Debes crear un nuevo tramo siempre que aparezca:
\begin{itemize}
  \item una carga puntual;
  \item un momento puntual;
  \item inicio o final de una carga repartida;
  \item cambio de expresión de $q(x)$;
  \item una articulación;
  \item un cambio de geometría.
\end{itemize}

\safesection{Paso 4: realiza el corte}
Coloca
\[
N(x),\qquad V(x),\qquad M(x)
\]
siguiendo siempre el mismo criterio de signos.

\safesection{Paso 5: comprueba con las relaciones diferenciales}
Una vez obtenidas las ecuaciones:
\[
N'=-q_x,\qquad V'=-q_y,\qquad M'=V.
\]
Si no se cumplen, existe casi con seguridad un error de signo o integración.

\safesection{Paso 6: comprueba puntos especiales}
Busca $V=0$ para localizar extremos de $M$.

Comprueba $M=0$ en articulaciones y extremos simplemente apoyados.

En un extremo libre sin carga puntual:
\[
V=M=0.
\]

\safesection{Paso 7: piensa físicamente}
Pregúntate:
\begin{itemize}
  \item ¿las reacciones tienen sentido?;
  \item ¿las unidades son correctas?;
  \item ¿el diagrama es continuo donde debe serlo?;
  \item ¿hay un salto donde existe una carga puntual?;
  \item ¿la curva tiene el grado esperado?
\end{itemize}

% ============================================================
\safechapter{Errores típicos que cuestan puntos}
% ============================================================

\safesection{Confundir carga distribuida con fuerza}
$q$ tiene unidades $\mathrm{kN/m}$, mientras que su resultante
\[
Q=\int q\dd x
\]
tiene unidades $\mathrm{kN}$.

\safesection{Usar mal el centroide de una carga triangular}
La resultante está a $L/3$ del lado de mayor intensidad, no del lado donde la carga vale cero.

\safesection{Confundir esfuerzo con tensión}
\[
N,V\rightarrow\text{fuerzas},
\qquad
M,T\rightarrow\text{momentos},
\qquad
\sigma,\tau\rightarrow\text{fuerza/área}.
\]

\safesection{No separar una estructura en una rótula}
En una rótula interna:
\[
M=0.
\]
Es una ecuación adicional.

\safesection{Cambiar el criterio de signos a mitad de ejercicio}
Un criterio distinto puede ser matemáticamente válido. Mezclar dos criterios no lo es.

\begin{tip}
Dibuja en una esquina del examen tu convenio positivo para $N$, $V$ y $M$ y no lo cambies.
\end{tip}

\safesection{Dibujar diagramas incompatibles con las derivadas}
Si $q=\text{cte.}$ no puedes dibujar $V=\text{cte.}$ porque $V'=-q\neq0$.

Y si $V$ es lineal, entonces $M$ es parabólico.

% ============================================================
